\documentclass[letterpaper]{article}

\usepackage[utf8]{inputenc}
\usepackage[spanish]{babel}

\usepackage{amssymb, amsmath}

\usepackage{geometry}
\geometry{margin=2cm}

\title{C\'alculo Integral}
\author{Jose Rosendo Camal Teh}
\date{15 de marzo de 2026}

\pagestyle{plain}

\begin{document}
   
	\begin{titlepage}
	    \centering
	    {\scshape\LARGE Universidad Aztlán \par}
	    \vspace{1cm}
	    {\scshape\Large Sede Playa del Carmen \par}
	    \vspace{3cm}
	    {\scshape\Huge Integrales \par}
	    \vspace{1cm}
	    {\itshape\Large Ejercicios\par}
	    \vspace{3cm}
	    {\Large José Rosendo Camal Teh \par}
	    \vfill
	    {\Large Ingeniería en Sistemas Computacionales \par}
	    {\Large Segundo Cuatrimestre \par}
	    {\Large C\'alculo Integral \par}
	    \vfill
	    {\Large Profe. Xicontecatl Carreto \par}
	    \vfill
	    {\Large 15 de marzo de 2026 \par}
	    \vfill
	    {\Large Playa del Carmen, Quintana Roo \par}
	\end{titlepage}

	
	\section*{C\'alculo Integral}
	\textbf{Instrucciones:} Calcular las siguientes integrales inmediatas.

	\subsection*{Ejercicio 1} % Primera integral inmediata
		\vspace{0.5cm}
		% Integral 1
		\[ \int \frac{9}{x^4} \, dx  \]
		\vspace{0.5cm}
		% Desarrollo
		\vspace{0.2cm}
		Desarollamos y buscamos la antiderivada:
		\begin{align*}
			\int \frac{9}{x^4} \, dx &= \int 9 \cdot \frac{1}{x^4} \, dx \\
			&= 9 \int \frac{1}{x^4} \, dx \\
			&= 9 \int x^{-4} \, dx \\
			&= 9 \left[ \frac{x^{-4 + 1}}{-4 + 1} \right] + C \\
			&= 9 \left[ \frac{x^{-3}}{-3} \right] + C \\
			&= 9 \cdot \frac{x^{-3}}{-3} + C \\
			&= 9 \cdot \left( \frac{1}{-3} \right) \cdot x^{-3} + C \\
			&= 9 \cdot \left( -\frac{1}{3} \right) \cdot x^{-3} + C \\
			&= -\frac{9}{3} x^{-3} + C \\
			&= -3 x^{-3} + C \\
			&= -\frac{3}{x^3} + C
		\end{align*}
		\vspace{0.2cm}
		Comprobamos:
		\begin{align*}
			\frac{9}{x^4} &= \frac{d}{dx} -\frac{3}{x^3} + C \\
			&= \frac{d}{dx} -\frac{3}{x^3} + \frac{d}{dx} C \\
			&= -3 \frac{d}{dx} \frac{1}{x^3} + 0 \\
			&= -3 \frac{d}{dx} x^{-3} \\
			&= -3 \left( -3 x^{-3 - 1} \right) \\
			&= -3 \left( -3 x^{-4} \right) \\
			&= 9 x^{-4} \\
			\frac{9}{x^4} &= \frac{9}{x^4}
		\end{align*}
		% Resultado
		\vspace{0.2cm}
		Resultado:
		\[ \int \frac{9}{x^4} \, dx = -\frac{3}{x^3} + C\]
		
	\subsection*{Ejercicio 2} % Segunda integral inmediata
		
		% Integral 2
		\vspace{0.5cm}
		\[ \int \frac{\, dx}{x^{-2}} \]
		\vspace{0.5cm}
		% Desarrollo
		\vspace{0.2cm}
		Desarrollamos la integral y buscamos la antiderivada:
		\begin{align*}
			\int \frac{\, dx}{x^{-2}} &= \int \frac{1}{x^{-2}} \, dx \\
			&= \int \frac{1}{ \frac{1}{x^2} } \, dx \\
			&= \int \frac{ \frac{1}{1} }{ \frac{1}{x^2} } \, dx \\
			&= \int \frac{x^2}{1} \, dx \\
			&= \int x^2 \, dx \\
			&= \frac{x^{2 + 1}}{2 + 1} + C \\
			&= \frac{x^3}{3} + C \\
			&= \frac{1}{3} x^3 + C
		\end{align*}
		\vspace{0.2cm}
		Comprobamos:
		\begin{align*}
			\frac{1}{x^{-2}} &= \frac{d}{dx} \frac{1}{3} x^3 + C \\
			&= \frac{d}{dx} \frac{1}{3} x^3 + \frac{d}{dx} C \\
			&= \frac{1}{3} \frac{d}{dx} x^3 + 0 \\
			&= \frac{1}{3} \left( 3 x^{3 - 1} \right) \\
			&= \frac{1}{3} \left( 3 x^2 \right) \\
			&= \frac{3}{3} x^2 \\
			\frac{1}{x^{-2}} &= x^2 \\
			\frac{1}{ \frac{1}{x^2}} &= \\
			x^2 &= x^2
		\end{align*}
		% Resultado
		\vspace{0.2cm}
		Resultado:
		\[ \int \frac{\, dx}{x^{-2}} = \frac{1}{3} x^3 + C \]
		
	\subsection*{Ejercicio 3} % Tercera integral inmediata
		
		% Integral 3
		\vspace{0.5cm}
		\[ \int x^{ \frac{2}{3} } \, dx \]
		\vspace{0.5cm}
		% Desarrollo
		\vspace{0.2cm}
		Desarrollamos la integral y buscamos la antiderivada:
		\begin{align*}
			\int x^{ \frac{2}{3} } \, dx &= \frac{x^{ \frac{2}{3} + 1}}{ \frac{2}{3} + 1} + C \\
			&= \frac{x^{ \frac{2}{3} + \frac{3}{3} }}{ \frac{2}{3} + \frac{3}{3} } + C \\
			&= \frac{x^{ \frac{2 + 3}{3} }}{ \frac{2 + 3}{3} } + C \\
			&= \frac{x^{ \frac{5}{3} }}{ \frac{5}{3} } + C \\
			&= \frac{ \sqrt[3]{x^5} }{ \frac{5}{3} } + C \\
			&= \frac{ \frac{ \sqrt[3]{x^5} }{1} }{ \frac{5}{3} } + C \\
			&= \frac{3  \sqrt[3]{x^5}}{5} + C \\
			&= \frac{3x \sqrt[3]{x^2}}{5} + C
		\end{align*}
		\vspace{0.2cm}
		Comprobamos:
		\begin{align*}
			x^{ \frac{2}{3} } &= \frac{d}{dx} \frac{3x \sqrt[3]{x^2}}{5} + C \\
			&= \frac{d}{dx} \frac{3x \sqrt[3]{x^2}}{5} + \frac{d}{dx} C \\
			&= \frac{3}{5} \frac{d}{dx} x \sqrt[3]{x^2} + 0 \\
			&= \frac{3}{5} \frac{d}{dx} \sqrt[3]{x^3} \sqrt[3]{x^2} \\
			&= \frac{3}{5} \frac{d}{dx} x^{ \frac{3}{3} } \cdot x^{ \frac{2}{3} } \\
			&= \frac{3}{5} \frac{d}{dx} x^{ \frac{3}{3} + \frac{2}{3} } \\
			&= \frac{3}{5} \frac{d}{dx} x^{ \frac{5}{3} } \\
			&= \frac{3}{5} \cdot \frac{5}{3} x^{\frac{5}{3} - 1} \\
			&= \frac{3}{5} \cdot \frac{5}{3} x^{ \frac{5}{3} - \frac{3}{3} } \\
			&= \frac{3}{5} \cdot \frac{5}{3} \cdot x^{ \frac{5 - 3}{3} } \\
			&= \frac{3}{5} \cdot \frac{5}{3} \cdot x^{ \frac{2}{3} } \\
			&= \left( \frac{3}{5} \cdot \frac{5}{3} \right) \cdot x^{ \frac{2}{3} } \\
			&= x^{ \frac{2}{3} } \\
			x^{ \frac{2}{3} } &= x^{ \frac{2}{3} }
		\end{align*}
		% Resultado
		\vspace{0.2cm}
		Resultado:
		\[ \int x^{ \frac{2}{3} } \, dx = \frac{ 3x \sqrt[3]{x^2}}{5} + C \]

\end{document}
